\section{卷积层详解}

\subsection{从全连接到卷积}

\begin{figure}[H]
\centering
\includegraphics[width=0.95\textwidth]{slides-images/slide-027.jpg}
\caption{全连接层 vs 卷积层：全连接层将整张图像展平后做内积；卷积层保持三维张量结构\protect\footnotemark}
\end{figure}
\footnotetext{Slides 第28页。}

全连接层中，输出向量的每个元素是权重矩阵的一行与\textbf{整个}输入向量的内积。这意味着：
\begin{itemize}
    \item 每个输出数值``看到''了图像中的\textbf{所有}像素
    \item 空间结构被完全丢弃
\end{itemize}

卷积层的核心改变：每个输出数值只依赖于输入图像的一个\textbf{局部区域}。

\subsection{卷积运算}

\begin{figure}[H]
\centering
\includegraphics[width=0.95\textwidth]{slides-images/slide-029.jpg}
\caption{单个卷积滤波器的工作方式：$5 \times 5 \times 3$ 的滤波器在 $32 \times 32 \times 3$ 的输入图像上滑动\protect\footnotemark}
\end{figure}
\footnotetext{Slides 第30页。}

卷积层的核心操作：
\begin{enumerate}
    \item 定义一组\textbf{滤波器}（filters / kernels），每个滤波器是一个小的三维张量，空间尺寸为 $K \times K$，深度等于输入通道数 $C_{in}$
    \item 将滤波器在输入图像上\textbf{滑动}，在每个位置计算滤波器与对应局部区域的\textbf{内积}（点积 + 偏置）
    \item 每个滤波器产生一个二维\textbf{激活图}（activation map）
    \item $C_{out}$ 个滤波器产生 $C_{out}$ 个激活图，堆叠后形成输出张量
\end{enumerate}

\begin{importantbox}{卷积层的参数与输出形状}
\textbf{输入}：$C_{in} \times H \times W$

\textbf{滤波器}：$C_{out}$ 个，每个大小为 $C_{in} \times K \times K$

\textbf{输出}：$C_{out} \times H' \times W'$

其中 $H' = \lfloor(H - K + 2P) / S\rfloor + 1$，$W'$ 类似。

\textbf{总参数量}：$C_{out} \times (C_{in} \times K \times K + 1)$（含偏置）
\end{importantbox}

\begin{figure}[H]
\centering
\includegraphics[width=0.95\textwidth]{slides-images/slide-031.jpg}
\caption{多滤波器卷积：6 个 $5 \times 5 \times 3$ 滤波器产生 $6 \times 28 \times 28$ 的输出\protect\footnotemark}
\end{figure}
\footnotetext{Slides 第32页。}

\subsection{参数 vs 超参数}

\begin{warningbox}{区分参数与超参数}
\begin{itemize}
    \item \textbf{超参数}（训练前设定，不可学习）：滤波器数量 $C_{out}$、滤波器大小 $K$、步幅 $S$、填充 $P$
    \item \textbf{参数}（通过梯度下降学习）：滤波器中每个元素的数值、偏置值
\end{itemize}
超参数决定了张量的\textbf{形状}，参数决定了张量的\textbf{数值}。滤波器在训练开始时随机初始化，然后通过反向传播不断优化。
\end{warningbox}

\subsection{卷积滤波器学到了什么？}

\begin{figure}[H]
\centering
\includegraphics[width=0.95\textwidth]{slides-images/slide-039.jpg}
\caption{AlexNet 第一层卷积滤波器的可视化：学到了颜色对比和边缘检测器\protect\footnotemark}
\end{figure}
\footnotetext{Slides 第40页。}

对第一层卷积滤波器的可视化揭示了两类模式：
\begin{itemize}
    \item \textbf{颜色检测器}：检测对比色（如红-绿对比、色块等）
    \item \textbf{边缘检测器}：检测水平、垂直、对角线方向的边缘
\end{itemize}

这个结果具有惊人的普遍性——几乎所有 CNN、所有数据集、所有任务中，第一层滤波器都学到类似的模式。

\begin{figure}[H]
\centering
\includegraphics[width=0.95\textwidth]{slides-images/slide-040.jpg}
\caption{深层卷积滤波器的可视化：更高层学到了更大的空间结构——纹理、物体部件、甚至完整物体\protect\footnotemark}
\end{figure}
\footnotetext{Slides 第41页。}

\begin{importantbox}{CNN 的层级特征学习}
CNN 自动学习了一个从低级到高级的\textbf{特征层次结构}：
\begin{itemize}
    \item 第 1 层：边缘和颜色
    \item 中间层：纹理和简单形状
    \item 高层：物体部件和语义概念
\end{itemize}
这与人类视觉系统中初级视皮层到高级视觉区域的层级处理有着惊人的相似性。
\end{importantbox}

\subsection{非线性激活函数不可或缺}

\begin{figure}[H]
\centering
\includegraphics[width=0.85\textwidth]{slides-images/slide-037.jpg}
\caption{堆叠卷积层但不加激活函数 = 一个卷积层（线性复合仍是线性）\protect\footnotemark}
\end{figure}
\footnotetext{Slides 第38页。}

卷积运算本质上是内积——一种线性运算。两个线性运算的复合仍然是线性运算。因此，如果直接堆叠多个卷积层而不加非线性激活函数，整个网络的表达能力等价于单个卷积层。

\begin{warningbox}{线性层堆叠 = 单层}
$f_2 \circ f_1(x) = W_2(W_1 x) = (W_2 W_1)x = W'x$

无论堆叠多少个线性层，等价于一个线性变换。\textbf{必须}在卷积层之间插入非线性激活函数（如 ReLU）才能增加网络的表达能力。
\end{warningbox}

\subsection{本章小结}

卷积层是 CNN 的核心组件。它通过局部连接和权值共享实现了两个关键优势：（1）保留图像的空间结构，（2）大幅减少参数量。滤波器通过梯度下降自动学习，无需人工设计。不同层的滤波器自然形成从边缘到物体的层级特征表示。非线性激活函数是多层卷积网络表达能力的关键保障。

%% ========== Part 5: 卷积运算细节 ==========

